% Revised 12 November 2026 (REVISION_BRIEF_2026-11-12.md): academic register, research questions, contributions.
\chapter{Introduction}
\label{ch:intro}

\section{Motivation}
\label{sec:intro-task}
\label{sec:intro-anny}

Human mesh recovery estimates the pose and shape parameters of a parametric body model for each person in an
image; the body model turns them into a triangle mesh of the body surface. This thesis addresses the multi-person, whole-body variant:
pose, shape and distance from the camera for every person, including people who overlap or are partly occluded.

Multi-HMR~\cite{baradel2024multihmr} solves this task in a single forward pass. A Vision Transformer (ViT) encodes the
image, a detection head selects one keypoint per person, and a Human Perception Head regresses each body and its
distance. The released Multi-HMR-Anny checkpoint (the released baseline) predicts bodies in the Anny body
model~\cite{bregier2025anny,naver2024multihmrcode}; multi-person Anny prediction is therefore not new, and this
checkpoint serves as the external baseline.

SMPL~\cite{loper2015smpl} represents shape by statistical coefficients learned from scans of adults, usually the
first \BgSmplBetas{}, which carry no direct semantic meaning. Anny instead uses \MethodNumPheno{} named phenotype
parameters, such as age and height, and \MethodNumBones{} bones for the whole body
(Section~\ref{sec:bg-bodymodels}, Table~\ref{tab:body-models}). Anny can therefore represent a child directly, which
does not imply that a network recognises a child in a photograph.

Two difficulties remain. The first is supervision: AnnyOne~\cite{bregier2025anny} provides synthetic images with
exact Anny labels but an appearance unlike real photographs, while 3DPW~\cite{vonmarcard2018_3dpw} provides real
video frames labelled in SMPL; no real images with Anny labels exist. The second is cost: large image encoders
have many parameters and long inference times. Knowledge distillation, in which a large
trained teacher guides the training of a smaller student, reduces the cost and keeps single-pass inference.

\section{Problem statement and research questions}
\label{sec:intro-requirements}

The input is an RGB image with its camera intrinsics (focal length and principal point), which define the viewing
ray of each pixel. For each detected person the model predicts an Anny body: bone rotations, phenotype values and the
camera-space position of one head joint, the anchor. The deployed model is a ViT-B student; its ViT-L teacher, trained
in this thesis, is the accuracy reference, and the released baseline the external reference. Four research
questions structure the thesis:
\begin{enumerate}
\item \textbf{Real-image supervision.} How can the SMPL labels of 3DPW be converted into Anny labels, and which
  aspects of pose and shape do they constrain?
\item \textbf{Compact reconstruction.} How much smaller, cheaper and faster is a ViT-B student distilled from a ViT-L
  teacher, how much accuracy does it lose against the teacher and the released baseline, and what does distillation
  contribute?
\item \textbf{Multi-person detection.} How does detection recall depend on the number of people in an image?
\item \textbf{Video.} Can the image-based student process a single-person video, and do inference-time temporal filtering and
  landmark alignment make its output steadier and better placed?
\end{enumerate}

The student processes one frame at a time; no multi-frame network is trained. Reconstruction is evaluated on a
held-out part of AnnyOne and on converted 3DPW, age and gender for the teacher on Relative Human photographs with
age-group and gender tags. No new shape labels for real images are collected, and age in years is not estimated.

\section{Contributions}
\label{sec:intro-contrib}

I implemented the Anny output head and its target construction on top of the
Multi-HMR code, the SMPL-to-Anny converter, the training and evaluation pipeline, the distillation set-up and the
video filter and alignment; the pretrained backbone weights come from the released checkpoints.
\begin{enumerate}
\item \textbf{Anny labels for real images.} A joint-based fitting pipeline converts the standard 3DPW splits into
  Anny labels by fitting corresponding joints and virtual torso joints; the labels were checked with mesh overlays
  and a joint comparison (Chapter~\ref{ch:data}).
\item \textbf{A compact reconstruction system.} A ViT-L teacher is trained on AnnyOne and converted 3DPW with two
  corrections to its supervision, consistent head-anchor targets and a fixed output rotation correction for flipped
  bones, both identified by diagnostic experiments (Appendix~\ref{ch:diagnostics}). A ViT-B student receives a copy
  of the teacher's prediction heads, which read its features through a learned projector, and imitates the teacher's
  features and outputs.
\item \textbf{An evaluation under one protocol.} All models are evaluated with the same labels and settings, with
  pose, placement and shape assessed separately, ablations of the teacher corrections, sensitivity
  analyses of training length and temperature, and detection recall by the number of people per image.
\item \textbf{A video pipeline.} The pose, shape and position of each tracked person are filtered over time, and
  position and depth are aligned to detected 2D landmarks, without retraining (Section~\ref{sec:method-temporal}). The student and the released baseline are compared on
  the same frames.
\end{enumerate}

\section{Thesis structure}
\label{sec:intro-outline}

Chapter~\ref{ch:background} introduces the body models and related work. Chapter~\ref{ch:data} describes the
datasets and the label conversion, on which the method of Chapter~\ref{ch:method} depends.
Chapter~\ref{ch:results} reports the experiments. Chapter~\ref{ch:discussion} answers the research questions and states the
limitations; Chapter~\ref{ch:conclusion} concludes. The appendices give reproduction details, the joint
correspondence, settings, additional figures, earlier models, the diagnostic experiments, the conversion history and
additional teacher results.
